\makesubsectionpage{Bewertung von Klassifikationsmodellen}

\begin{frame}{Bewertung von Klassifikationsmodellen}
\begin{columns}
\column{0.60\textwidth}
Die Güte des trainierten Modells haben wir bisher mit \texttt{model.score()}
bewertet.

\vspace{0.4em}
\begin{itemize}
  \item \texttt{model.score()} liefert standardmäßig die
        \textbf{Trefferquote}.
  \item Also: korrekt gelabelte Testdaten geteilt durch die Gesamtzahl der
        Testdaten.
\end{itemize}

\vspace{0.4em}
\begin{block}{Aber}
In vielen Fällen ist die reine Trefferquote \textbf{kein geeignetes
Gütemaß}.
\end{block}

\column{0.35\textwidth}
\centering
\twemoji[width=0.8\linewidth]{bar chart}
\end{columns}
\end{frame}

\begin{frame}{Bewertung: Unbalancierte Datensätze}
\begin{columns}
\column{0.58\textwidth}
Datensätze mit ungleicher Klassenverteilung nennt man
\textbf{unbalanciert}. Reale Datensätze sind oft stark unbalanciert.

\vspace{0.4em}
\textbf{Beispiel:} Betrugserkennung bei Kreditkarten-Transaktionen
\begin{itemize}
  \item \textbf{99,98\,\%} zulässig, nur \textbf{0,02\,\%} Betrug.
  \item Ein Klassifikator, der \textbf{alles erlaubt}, erreicht
        99,98\,\% Trefferquote.
\end{itemize}

\vspace{0.4em}
\begin{blockC2}{Trugschluss}
Hohe Trefferquote -- und trotzdem \textbf{völlig nutzlos}: kein einziger
Betrug wird erkannt.
\end{blockC2}

\column{0.38\textwidth}
\centering
\resizebox{0.9\linewidth}{!}{%
\begin{tikzpicture}[
  kugel/.style={circle, thick, minimum size=10pt, inner sep=0pt}
]
\draw[very thick] (-0.1,2.0) -- (-0.1,0) -- (2.1,0) -- (2.1,2.0);
% viele "zulaessig" (blau), ein "Betrug" (rot)
\foreach \x/\y in {0.5/0.4, 1.0/0.4, 1.5/0.4, 0.5/0.9, 1.5/0.9,
                   0.5/1.4, 1.0/1.4, 1.5/1.4} {
  \node[kugel, fill=PrimaryColor, draw=PrimaryTitle] at (\x,\y) {};
}
\node[kugel, fill=AccentColor, draw=AccentTitle] at (1.0,0.9) {};
\end{tikzpicture}}\\[4pt]
{\footnotesize\itshape 1 Betrug unter vielen zulässigen Fällen.}
\end{columns}
\end{frame}



\begin{frame}{Bewertung: Konfusionsmatrix}
\begin{columns}
\column{0.4\textwidth}
Eine \textbf{Konfusionsmatrix} (confusion matrix) veranschaulicht, wie
viele Datensätze der Klassifikator richtig und falsch zugeordnet hat.

\vspace{0.4em}
Dafür brauchen wir:
\begin{itemize}
  \item die \textbf{Ergebnisse} des Klassifikators für alle Testdaten,
  \item die \textbf{erwarteten} Ergebnisse -- also die Labels der
        Testdatenmenge.
\end{itemize}

\column{0.55\textwidth}
\centering
\resizebox{\linewidth}{!}{%
\begin{tikzpicture}[
  ekopf/.style={fill=ExampleTitle, text=white, font=\footnotesize\bfseries,
                minimum width=3.4cm, minimum height=1.1cm},
  skopf/.style={fill=ExampleTitle, text=white, font=\footnotesize\bfseries,
                minimum width=2.6cm, minimum height=1.1cm},
  zelle/.style={font=\footnotesize, align=center, text width=2.4cm,
                minimum width=2.6cm, minimum height=1.1cm}
]
% Titel oben
\node[font=\footnotesize\bfseries] at (3.5,2.0) {Tatsächlicher Zielwert};

% Kopfzeile (Spaltenköpfe)
\node[ekopf] at (1.0,1.25) {};
\node[skopf] at (3.5,1.25) {POSITIV};
\node[skopf] at (6.1,1.25) {NEGATIV};

% erste Datenzeile (kleiner Abstand zur Kopfzeile)
\node[ekopf] at (1.0,0.1)  {POSITIV};
\node[zelle, fill=ExampleTitle!18] at (3.5,0.1) {Anzahl \textbf{TP}\\(true positive)};
\node[zelle, fill=ExampleTitle!8]  at (6.1,0.1) {Anzahl \textbf{FP}\\(false positive)};

% zweite Datenzeile
\node[ekopf] at (1.0,-1.05)  {NEGATIV};
\node[zelle, fill=ExampleTitle!8]  at (3.5,-1.05) {Anzahl \textbf{FN}\\(false negative)};
\node[zelle, fill=ExampleTitle!18] at (6.1,-1.05) {Anzahl \textbf{TN}\\(true negative)};

% Achsenbeschriftung links
\node[font=\footnotesize\bfseries, rotate=90] at (-0.9,-0.5)
  {Vorhersage};
\end{tikzpicture}}
\end{columns}
\end{frame}

\begin{frame}{Konfusionsmatrix: Beispiel Spritzguss}
Beim \textbf{Spritzgieß-Beispiel} ist das Label, ob ein Zyklus ein
\textbf{Gutteil} (positiv) oder ein \textbf{Schlechtteil} (negativ)
erzeugt hat. Die vier Möglichkeiten:

\vspace{0.6em}
\begin{columns}[T]
\column{0.48\textwidth}
\begin{blockC1}{Korrekt klassifiziert}
\begin{itemize}
  \item \textbf{TP} (true positive): Klassifikator sagt Gutteil -- war ein
        Gutteil.
  \item \textbf{TN} (true negative): Klassifikator sagt Schlechtteil --
        war ein Schlechtteil.
\end{itemize}
\end{blockC1}

\column{0.47\textwidth}
\begin{blockC2}{Falsch klassifiziert}
\begin{itemize}
  \item \textbf{FP} (false positive): Klassifikator sagt Gutteil -- war ein
        Schlechtteil.
  \item \textbf{FN} (false negative): Klassifikator sagt Schlechtteil --
        war ein Gutteil.
\end{itemize}
\end{blockC2}
\end{columns}
\end{frame}

\begin{frame}{Konfusionsmatrix: konkretes Beispiel}
\begin{columns}
\column{0.47\textwidth}
Wir klassifizieren \textbf{200 Produktionszyklen}. Die Konfusionsmatrix
könnte wie rechts aussehen.

\vspace{0.4em}
\begin{itemize}
  \item Sie zeigt, \textbf{welche Art} von Fehlern der Klassifikator macht.
  \item Daraus lassen sich leicht \textbf{alternative Gütemaße} berechnen.
\end{itemize}

\column{0.48\textwidth}
\centering
\resizebox{\linewidth}{!}{%
\begin{tikzpicture}[
  ekopf/.style={fill=ExampleTitle, text=white, font=\footnotesize\bfseries,
                minimum width=3.2cm, minimum height=1.2cm},
  skopf/.style={fill=ExampleTitle, text=white, font=\footnotesize\bfseries,
                minimum width=2.6cm, minimum height=1.2cm},
  zelle/.style={font=\large, align=center, text width=2.4cm,
                minimum width=2.6cm, minimum height=1.2cm}
]
% Titel oben
\node[font=\footnotesize\bfseries] at (3.4,2.2) {Tatsächlicher Zielwert};

% Kopfzeile
\node[ekopf] at (0.9,1.3) {};
\node[skopf] at (3.4,1.3) {POSITIV};
\node[skopf] at (6.0,1.3) {NEGATIV};

% erste Datenzeile
\node[ekopf] at (0.9,0.0)  {POSITIV};
\node[zelle, fill=ExampleTitle!18] at (3.4,0.0) {120};
\node[zelle, fill=ExampleTitle!8]  at (6.0,0.0) {20};

% zweite Datenzeile
\node[ekopf] at (0.9,-1.3)  {NEGATIV};
\node[zelle, fill=ExampleTitle!8]  at (3.4,-1.3) {40};
\node[zelle, fill=ExampleTitle!18] at (6.0,-1.3) {20};

% Achsenbeschriftung links
\node[font=\footnotesize\bfseries, rotate=90] at (-0.9,-0.7) {Vorhersage};
\end{tikzpicture}}
\end{columns}
\end{frame}


\begin{frame}{Bewertung: Sensitivität (Recall)}
\begin{columns}
\column{0.55\textwidth}
Die \textbf{Sensitivität} (True-Positive-Rate, Recall) gibt an, wie viele
der positiven Datenpunkte auch tatsächlich als positiv erkannt werden.

\vspace{0.4em}
\[
  \text{Recall} = \frac{\text{TP}}{\text{TP} + \text{FN}}
\]

\vspace{0.4em}
Sie ist wichtig, wenn ein Modell \textbf{alle positiven Fälle sicher
entdecken} soll.

\column{0.42\textwidth}
\begin{blockC1}{Anwendung: Krebserkennung}
Ein Klassifikator soll Krebszellen erkennen.
\begin{itemize}
  \item \textbf{Alle} tatsächlichen Krebsfälle müssen erkannt werden
        (frühe Behandlung).
  \item Ein \textbf{Fehlalarm} ist unkritisch -- eine zusätzliche
        Untersuchung schadet nicht.
\end{itemize}
\end{blockC1}
\end{columns}
\source{\cite{databasecamp2022konfusionsmatrix}}
\end{frame}


\begin{frame}{Bewertung: Spezifität}
\begin{columns}
\column{0.55\textwidth}
Die \textbf{Spezifität} (True-Negative-Rate) misst, wie viele der als
negativ klassifizierten Datenpunkte auch tatsächlich negativ sind.

\vspace{0.4em}
\[
  \text{Spezifität} = \frac{\text{TN}}{\text{TN} + \text{FP}}
\]

\vspace{0.4em}
Sie ist wichtig, wenn \textbf{Fehlalarme teuer} sind.

\column{0.42\textwidth}
\begin{blockC1}{Anwendung: Qualitätskontrolle}
Überwachung eines \textbf{kostenintensiven} Produkts.
\begin{itemize}
  \item Fehlerhafte Produkte sollen nicht zu den Kund*innen gelangen.
  \item Zugleich möglichst \textbf{wenige} Produkte unnötig vernichten.
\end{itemize}
\end{blockC1}
\end{columns}
\source{\cite{databasecamp2022konfusionsmatrix}}
\end{frame}

\begin{frame}{Bewertung: Trefferquote und Fehlerrate}
\begin{columns}
\column{0.49\textwidth}
\begin{blockC1}{Trefferquote (Accuracy)}
Anteil der korrekt klassifizierten Datensätze:
\[
  \frac{\text{TP} + \text{TN}}
       {\text{TP} + \text{TN} + \text{FP} + \text{FN}}
\]
\end{blockC1}

\column{0.48\textwidth}
\begin{blockC2}{Fehlerrate}
Anteil der falsch klassifizierten Datensätze:
\[
  \frac{\text{FP} + \text{FN}}
       {\text{TP} + \text{TN} + \text{FP} + \text{FN}}
\]
\end{blockC2}
\end{columns}

\vspace{0.6em}
\centering
{\footnotesize\itshape Beide ergänzen sich zu 1: Trefferquote $+$
Fehlerrate $= 1$.}
\source{\cite{databasecamp2022konfusionsmatrix}}
\end{frame}


\begin{frame}{Bewertung: Balancierte Fehlerrate (BER)}
\begin{columns}
\column{0.55\textwidth}
Bei stark \textbf{unbalancierten} Datensätzen verwendet man häufig die
\textbf{Balancierte Fehlerrate} (BER).

\vspace{0.4em}
\begin{itemize}
  \item Die Fehlerrate wird für jede Klasse \textbf{separat} berechnet.
  \item Die BER ist das \textbf{arithmetische Mittel} dieser einzelnen
        Fehlerraten.
\end{itemize}

\vspace{0.4em}
Man bestimmt also, wie viel Prozent der tatsächlich \textbf{positiven} und
wie viel Prozent der tatsächlich \textbf{negativen} Datensätze
\textbf{falsch} klassifiziert werden -- und mittelt beide Werte.

\column{0.4\textwidth}
\begin{blockC1}{Formel}
\[
  \text{BER} = \frac{1}{2}\!\left(
    \frac{\text{FN}}{\text{TP} + \text{FN}}
    +
    \frac{\text{FP}}{\text{TN} + \text{FP}}
  \right)
\]
\vspace{0.2em}

{\footnotesize Mittel aus Fehlerrate der positiven und der negativen
Klasse.}
\end{blockC1}
\end{columns}
\end{frame}

\begin{frame}{BER am Beispiel: Kreditkarten-Betrug}
\begin{columns}
\column{0.48\textwidth}
\textbf{10.000 Transaktionen}, davon nur 20 Betrugsfälle (positiv).
Der Klassifikator erkennt nur die Hälfte der Betrugsfälle:

\vspace{0.4em}
\centering
\resizebox{\linewidth}{!}{%
\begin{tikzpicture}[
  ekopf/.style={fill=ExampleTitle, text=white, font=\scriptsize\bfseries,
                minimum width=2.0cm, minimum height=1.0cm},
  skopf/.style={fill=ExampleTitle, text=white, font=\scriptsize\bfseries,
                minimum width=1.8cm, minimum height=1.0cm},
  zelle/.style={font=\normalsize, align=center, text width=1.7cm,
                minimum width=1.8cm, minimum height=1.0cm}
]
\node[font=\scriptsize\bfseries] at (2.5,1.9) {Tatsächlich};
\node[ekopf] at (0.7,1.2) {};
\node[skopf] at (2.5,1.2) {Betrug};
\node[skopf] at (4.3,1.2) {zulässig};
\node[ekopf] at (0.7,0.1)  {Betrug};
\node[zelle, fill=ExampleTitle!18] at (2.5,0.1) {10};
\node[zelle, fill=ExampleTitle!8]  at (4.3,0.1) {10};
\node[ekopf] at (0.7,-1.0)  {zulässig};
\node[zelle, fill=ExampleTitle!8]  at (2.5,-1.0) {10};
\node[zelle, fill=ExampleTitle!18] at (4.3,-1.0) {9970};
\end{tikzpicture}}

\column{0.47\textwidth}
\begin{block}{Trefferquote}
\[
  \frac{10 + 9970}{10000} = \mathbf{99{,}8\,\%}
\]
Sieht hervorragend aus.
\end{block}

\begin{blockC2}{Balancierte Fehlerrate}
\[
  \text{BER} = \frac{1}{2}\!\left(
    \underbrace{\tfrac{10}{20}}_{50\%}
    +
    \underbrace{\tfrac{10}{9980}}_{0{,}1\%}
  \right) \approx \mathbf{25{,}0\,\%}
\]
Entlarvt das schlechte Modell.
\end{blockC2}
\end{columns}
\end{frame}

\makequizpage{\QUIZURL}{\QUIZQRCODE}{\QUIZIMAGE}


\begin{frame}[fragile]{Konfusionsmatrix in Python}
\begin{columns}
\column{0.58\textwidth}
\begin{minted}[fontsize=\footnotesize]{python}
from sklearn import metrics

vorhersagen = model.predict(X_test)
confusion_matrix = metrics.confusion_matrix(
    y_test, vorhersagen)
print(confusion_matrix)
\end{minted}

\vspace{0.4em}
\begin{minted}[fontsize=\footnotesize, frame=none]{text}
[[31  2]
 [ 5 37]]
\end{minted}

\column{0.35\textwidth}
Die Konfusionsmatrix zeigt, welche Klassen der Klassifikator
\textbf{verwechselt}.

\vspace{0.4em}
\begin{itemize}
  \item \texttt{confusion\_matrix} wird aus den \textbf{vorhergesagten}
        Labels und den \textbf{Test}-Labels erstellt.
  \item Ein Aufruf liefert alle Werte auf einmal.
\end{itemize}
\end{columns}
\source{Entscheidungsbaum Confusion Matrix.ipynb}
\end{frame}

\begin{frame}[fragile]{Konfusionsmatrix visualisieren}
\begin{columns}
\column{0.52\textwidth}
Mit \textbf{matplotlib} lässt sich die Konfusionsmatrix anschaulicher
darstellen -- inklusive benannter Label.

\vspace{0.4em}
\begin{minted}[fontsize=\footnotesize]{python}
import matplotlib.pyplot as plt
import numpy
from sklearn import metrics

cm_display = 
    metrics.ConfusionMatrixDisplay(
        confusion_matrix=confusion_matrix,
        display_labels=["gut", "schlecht"])
cm_display.plot()
plt.show()
\end{minted}

\column{0.43\textwidth}
\centering
\includesvg[width=\linewidth]{figures/confusion_matrix_baum.svg}
\end{columns}
\source{Entscheidungsbaum Confusion Matrix.ipynb}
\end{frame}

\begin{frame}{Übung: Konfusionsmatrix für Iris}
\begin{columns}

\column{0.65\textwidth}
Erweitern Sie Ihr Notebook mit dem \textbf{Iris-Datensatz} so, dass die
\textbf{Konfusionsmatrix} Ihres Klassifikators berechnet und ausgegeben
wird.

\vspace{0.4em}
\begin{itemize}
  \item Hat die Matrix die \textbf{Form}, die Sie erwartet haben?
  \item Welche Unterart wird mit welcher anderen \textbf{verwechselt}?
\end{itemize}

\column{0.30\textwidth}
\centering
\twemoji[width=0.8\linewidth]{laptop}


\end{columns}
\source{Lösung: Entscheidungsbaum Konfusionsmatrix.ipynb}
\end{frame}



\makequestionspage{\FRAGENIMAGE}